This work presents the modeling of a system capable of detect and track objects for several applications in Surveillance and Monitoring. Objective, motivation, literature review, and the organization of this work are presented in the following sections.

\section{Objective}
In the context of real-time visual object tracking systems, a persistent challenge lies in achieving reliable performance, especially when considering the complex dynamics of the actuator and the communication system. The accuracy of tracking in images, acquired and processed in real-time, critically depends on the ability to mitigate inherent nonlinearities in mechatronic systems and to compensate for communication delays, ensuring that the object of interest remains centered within the field of view. In this scenario, this project aims to develop a comprehensive control solution for object tracking applications in images.

\section{Motivation}
In recent years, technology has transformed many sectors, which are now significantly benefiting from systems capable of detecting and tracking objects. These systems prove invaluable in diverse situations, ranging from enhancing security applications to optimizing industrial production, enabling the advancement of autonomous vehicles and empowering robots to perform complex tasks. The ability to perceive and track objects in real-time not only boosts efficiency but also considerably improves the safety and accuracy of numerous daily activities and critical operations. In public safety, for instance, cameras equipped with these technologies are indispensable for monitoring busy areas, deterring criminal activities, and assisting in investigations. They can identify suspicious behavior and instantly alert law enforcement, providing detailed imagery as evidence. This marks a substantial progression from traditional surveillance methods, which are heavily reliant on human operators. Within the automotive sector, autonomous vehicles depend on these systems to navigate securely, accurately identifying pedestrians, other vehicles, and obstacles on the road. Furthermore, these systems empower vehicles to make swift and precise decisions, adapting seamlessly to environmental changes. As we move closer to a future characterized by increased vehicle autonomy, the imperative to continuously enhance these systems becomes ever more evident. In military operations, object detection and tracking are paramount for surveillance and security. Drones outfitted with these technologies monitor conflict-affected zones, offering a detailed, real-time perspective of the theater of operations. These systems are vital for identifying and tracking threats, such as adversarial vehicles and equipment, thereby augmenting operational effectiveness and mitigating risks for personnel. The speed and precision in target detection are critical determinants of success in military missions, and advancements in this domain can profoundly impact critical situations.

Integrating object detection and tracking technologies with multi-axis gimbal systems elevates the performance of these applications to a new level. Gimbals provide a mechanical platform capable of isolating a sensor from undesirable base movements (image stabilization) and actively reorienting the Line-of-Sight (LOS) to follow a target (precise tracking). This integration not only yields more stable and smoother footage, crucial for the quality of visual information, but also drastically extends the sensor's operational field of view, enabling the monitoring of distant or fast-moving targets across expansive areas.

In the aeronautical sector, the application of gimbals equipped with optoelectronic sensors is particularly pertinent. Such solutions form the core of advanced Intelligence, Surveillance, and Reconnaissance (ISR) systems utilized in Unmanned Aerial Vehicles (UAVs). In these operations, the ability to maintain a target of interest at the center of the camera's field of view, regardless of the aircraft's or the target's own movements, is crucial for collecting high-quality data and ensuring mission effectiveness.

Despite significant progress in the field, developing robust and high-performance visual tracking systems still presents considerable challenges. The inherent complexity of gimbal dynamics, which involves nonlinearities (such as friction) in addition to communication delays intrinsic to distributed system architectures, poses an obstacle to achieving precise and reliable tracking. The balance between object detection accuracy and the computational cost required for real-time processing adds another layer of complexity to this challenge.

Within this context, this Master's thesis aims to develop a control solution for gimbal-based object tracking systems using real-time acquired images. The motivation stems from the existing gap between simplified theoretical models and the actual behavior of complex mechatronic systems. By modeling the dynamics and nonlinearities of the gimbal and its actuators, combined with the design of a controller capable of mitigating the effects of disturbances and delays, this work seeks to contribute to improving the performance of visual tracking systems in engineering applications.

\section{Contribution of the Work}

While adopting strategic simplifications in its modeling to maintain clarity and project feasibility, this work offers relevant academic and practical contributions in the field of visual object tracking through gimbal systems. 

The study proposes an integrated and complete system architecture, uniting visual perception (YOLOv8 algorithm), data communication (UDP protocol), and motion control (PID controller) with the electromechanical dynamics of a two-degrees-of-freedom gimbal. This end-to-end integration, simulated in a unified environment, represents a fundamental step for autonomous tracking systems. The system modeling, despite the inherent simplicity of the linearized model for controller design, explicitly and modularly incorporates the most critical nonlinearities for gimbal behavior, such as Coulomb friction, as well as the inclusion of a communication delay via Padé approximation, which adds an important level of realism. Validation is performed in a Simulink environment, with a practical focus on real-time, allowing for a more faithful evaluation of controller performance. Additionally, the developed model establishes a clear platform for future research and for the introduction of additional complexities, by isolating and modeling primary nonlinearities as disturbances, offering a didactic starting point for the development of more advanced control techniques or the incorporation of further details.

\section{Literature Review}
A methodology for detecting and tracking objects in images is proposed by \citen{Mane}. The approach consists of using a CNN (convolutional neural network) to perform two procedures, object detection and tracking. In this work, the TF (TensorFlow) tool is presented to simplify the construction of the object detection algorithm.

\citen{6909475} present a system that perform object detection in three steps. The first generates category-independent region proposal, and these proposals define the set of candidate detections available. The second module is a large convolutional neural network that extracts a fixed-length feature vector from each region. The third module is a set of class specific linear SVMs. This system is called R-CNN (\textit{Region-based Convolutional Neural Networks}). 

R-CNN achieves excellent object detection accuracy according to \citen{7410526}, however, it has some drawbacks, such as expensive training and slow object detection. By making changes on the training and test methods, \citen{7410526} presented a new training algorithm, called Fast R-CNN. As the name suggests, Fast R-CNN is faster to train and test in relation to R-CNN.

\citen{7485869} presents a different region proposal method, called Faster R-CNN. Faster R-CNN is composed of two modules, where the first is a deep fully convolutional network that proposes regions, and the second is a Fast R-CNN that uses the proposed regions. This new approach can achieve an speed of 5-17 \textit{frames per second}, which is significantly higher than other region proposal networks.

\citen{YOLO} proposes a solution for detecting objects in videos, called ``You Only Look Once'' (YOLO). This methodology differs because it unifies different components of object detection in a single neural network, thus managing to use information from the image as a whole to predict each bounding box (region of the image that locates the detected object). The YOLO model is implemented using a CNN.

An application of YOLO in its most current version (YOLOv8) to solve a surveillance problem, by detecting and tracking movement in real time, is presented in \citen{YOLOsurveillance}. The approach was to use a dataset for security applications and augment the YOLOv8 architecture with security-specific algorithms to improve real-time threat detection.

SSD (Single Shot Detector) algorithm is proposed by \citen{Liu_2016}. According to \citen{Liu_2016}, algorithms such as Faster R-CNN are accurate, however, for embedded systems is computationally intensive, and the cost of increasing speed in these approaches comes with significantly decrease in detection acuracy. The results on the Pascal VOC2007 test shows that SSD can achieve similar accuracy to Faster R-CNN, with higher speed.

The work done by \citen{Kalmanxalpha} presents a comparison between two algorithms for solving the problem of tracking objects in images, the Kalman filter and the alpha-beta filter. Using a sphere as the object of study, performance is given by analyzing parameters such as the estimate of the centroid, the radius of the sphere estimated by the algorithm and the iteration time. The comparison between the algorithms shows greater accuracy in the results presented by the use of the Kalman filter. 

In \citen{GimbalYOLO}, an airborne inertial platform (\textit{``Gimbal''}) is controlled based on images. In order to achieve real-time performance, the author points to the use of algorithms such as YOLO and SSD, which can be called ``one-stage detectors'' and can speed up the process of locating objects in an image.
 
\citen{Rahmat2016Sliding} conducted a comprehensive study on a two-degrees-of-freedom (2-DOF) gun turret model, providing a detailed derivation of its dynamics and simulating its behavior in MATLAB-Simulink. Their work specifically evaluated and compared the performance of PID and Sliding Mode Control (SMC) strategies for accurate target angle tracking, demonstrating the effectiveness of the SMC controller in generating sufficient drive torque and precisely reducing system error.

\citen{Galal2009Fuzzy} propose a solution for the control of a gun-turret system, a complex multi-body, multi-input, multi-output (MIMO) nonlinear system, utilizing Fuzzy Logic Control (FLC). This methodology addresses the challenges of achieving rapid and precise tracking responses in weapon systems, particularly under the influence of disturbances, nonlinearities, and modeling uncertainties that often degrade the performance of conventional controllers. The work details the application of a fuzzy scheme for compensation and nonlinear feedback control laws, which are crucial for managing the coupled dynamics of the turret and barrel driven by two DC motors. The authors highlight that FLC provides an effective means of capturing the approximate, inexact nature of the real world and adapting to unexpected parameter variations, thereby offering a robust and adaptive control strategy that is designed to accommodate inherent system uncertainties. The effectiveness of their fuzzy logic controller design is validated through simulation results, demonstrating improved system tracking performance.

\citen{TONG2020105312} investigated the gimbal torque and coupling torque in a six-degrees-of-freedom (6-DOF) magnetically suspended yaw gimbal, exploring various control methods such as PID, fuzzy control, Sliding Mode Control (SMC), adaptive control, feed-forward control, disturbance observers, and H-infinity control, and specifically employing decentralized PD control for radial active magnetic bearings. Their research focused on developing dynamic models of the yaw gimbal and a three-axis inertially stabilized platform (ISP), analyzing coupling torques among the gimbals, and studying the characteristics of gimbal torque generated by the magnetic suspension system. The experimental results indicated that the magnetic suspension system's gimbal torque can effectively compensate for coupling torques, thereby improving the yaw gimbal's attitude stabilization precision.

\citen{poyrazoglu_2017} presented a detailed modeling and control strategy for a two-degrees-of-freedom (2-DOF) gimbal system, employing the Newton-Euler approach to incorporate significant nonlinearities such as friction and mass unbalance. The study further explored state-space representations for both coupled and decoupled gimbal models, designing and comparing cascade PI and Linear Quadratic Integral (LQI) controllers for stabilization and target tracking within a MATLAB/Simulink virtual environment.

\citen{espinosa2016diseno} presented a work on the design and implementation of a control system for a two-degrees-of-freedom (2-DOF) gimbal, intended for UAV intelligence, surveillance, and reconnaissance missions. His work focused on developing a control structure to stabilize the camera's optical axis against environmental disturbances, deriving equations of motion using the Newton-Euler method, and validating Sliding Mode Control (SMC) algorithms through both simulation and experimental tests on a physical prototype, comparing their performance against a Proportional-Integral (PI) compensator.

\section{Organization}
This work is organized into 7 chapters, following a logical progression that spans from theoretical fundamentals and detailed system modeling to controller design, its validation through simulation, and finally, conclusions and suggestions for future work.

Chapter 1 establishes the general context of the research followed by the motivation, definition of the problem to be investigated, and the objectives and expected contributions of this study, whilst presenting an overview of existing literature.

Chapter 2 addresses the fundamental concepts of computer vision, neural networks, gimbals and stabilized platforms, vision-based object tracking, and control of pointing systems, in addition to analyzing related works that support the development of this project.

Chapter 3 is dedicated to the construction of the complete system's mathematical model, detailing the gimbal's kinematics and dynamics, the DC motor's modeling with its essential nonlinearities, the sensors used for feedback, and the representation of communication delay.

Chapter 4 describes the methodology used to design and tune the controller, specifying the chosen control strategy, the process of preparing the visual tracking input signal, and the performance analysis of the controller.

Chapter 5 presents the simulation environment and the obtained results, validating the proposed control system's performance.

Finally, Chapter 6: Results and Discussion (or Chapter 7: Conclusions and Future Work, depending on the final numbering) consolidates the main conclusions of the work, discusses its limitations, and proposes directions for future research in the field.